\documentclass{article}
\usepackage{amsmath, amsfonts, amssymb, amsthm, mathtools, mathrsfs}
\usepackage[usenames,dvipsnames]{xcolor}
\usepackage[margin=1in]{geometry}
\usepackage{fancyhdr}
\pagestyle{fancy}
%matrix
\makeatletter
\renewcommand*\env@matrix[1][c]{\hskip -\arraycolsep
  \let\@ifnextchar\new@ifnextchar
  \array{*\c@MaxMatrixCols #1}}
\makeatother
\usepackage{thmtools}
\usepackage{tikz}
\usepackage{tikz-cd}
\usepackage[framemethod=TikZ]{mdframed}
\mdfsetup{skipabove=1em,skipbelow=0em, innertopmargin=5pt, innerbottommargin=6pt}
\theoremstyle{definition}
\declaretheoremstyle[headfont=\bfseries\sffamily, bodyfont=\normalfont, mdframed={nobreak}]{thmbox}
\declaretheoremstyle[headfont=\bfseries\sffamily, bodyfont=\normalfont, mdframed={bottomline=false, topline=false, rightline=false}, qed=\(\blacksquare\)]{proofline}
\declaretheoremstyle[headfont=\bfseries\sffamily, bodyfont=\normalfont, qed=\(\blacksquare\)]{soln}
\declaretheorem[numbered=yes, style=thmbox, name=Exercise]{exercise}
\declaretheorem[numbered=no, style=thmbox, name=Theorem]{theorem}
\declaretheorem[numbered=no, style=thmbox, name=Lemma]{lemma}
\declaretheorem[numbered=no, style=proofline, name=Proof]{replacementproof}
\declaretheorem[numbered=no, style=soln, name=Solution]{replacementsoln}
\renewenvironment{proof}[1][\proofname]{\begin{replacementproof}}{\end{replacementproof}}
\newenvironment{solution}[1][]{\begin{replacementsoln}}{\end{replacementsoln}}
\begin{document}
\thispagestyle{fancy}
\pagestyle{fancy}
\fancyhead[R]{Grant Talbert}
\fancyhead[L]{MAT 202}
\begin{center}
\textbf{Written Assignment 6}\\
Grant Talbert
\end{center}
\begin{exercise}
  Let $B = \left\{(1, 0), (0, 1)\right\}$ and $B' = \left\{(1, 1), (1, 2)\right\}$ be bases for $\mathbb{R}^2$.
 \begin{enumerate}
 \item Find the matrix $A$ for $T:\mathbb{R}^2\rightarrow\mathbb{R}^2$, $T(x, y) = (x-2y, x+4y)$, relative to the basis $B$.
 \item Find the transition matrix from $\mathcal{B} $ to $\mathcal{B}^{\prime} $.
 \item Find the matrix $A'$ of $T$ relative to the basis $B'$.
 \item Find $[T(\vec{v})]_{B'}$ when $[\vec{v}]_{B'} = \begin{bmatrix*}[r]
 3\\-2
 \end{bmatrix*}$.
 \end{enumerate}
\end{exercise}
\begin{solution}
  Trivially, \(T(1,0)=(1,1)\) and \(T(0,1)=(-2,4)\). Define
  \[
    A \coloneqq \begin{bmatrix}[r]
      1 & -2  \\
      1 &4   \\
    \end{bmatrix}
  \]
  and notice that \(T\) has the representation \(T:\mathbf{x}\mapsto A \mathbf{x}\).
  \[
    \begin{bmatrix}[r]
      1 &1  &1  &0   \\
       1&2  &0  &1   \\
    \end{bmatrix}\xlongrightarrow{rref} \begin{bmatrix}[r]
      1 &0  &2  &-1   \\
       0&1  &-1  &1   \\
    \end{bmatrix}
  \]
  Thus, the matrix
  \[
    P ^{-1} =\begin{bmatrix}[r]
      2 &-1   \\
       -1&   1\\
    \end{bmatrix}
  \]
  gives a transition from the basis \(\mathcal{B} \) to \(\mathcal{B} ^{\prime} \). Also notice that its inverse will be
  \[
    P=\begin{bmatrix}[r]
      1 &1   \\
       1&2   \\
    \end{bmatrix}
  \]
  which gives a transition from \(\mathcal{B} ^{\prime} \) to \(\mathcal{B} \). Thus, we also know that \(A^{\prime} \) is given as
  \[
    A^{\prime} =P ^{-1} AP = \begin{bmatrix}[r]
      2 &-1   \\
       -1&   1\\
    \end{bmatrix}\begin{bmatrix}[r]
      1 & -2  \\
      1 &4   \\
    \end{bmatrix}\begin{bmatrix}[r]
      1 &1   \\
       1&2   \\
    \end{bmatrix} = \begin{bmatrix}[r]
      -7 &-15   \\
       6&12   \\
    \end{bmatrix}
  \]
  Thus we have
  \[
    [T(\mathbf{v})]_{\mathcal{B} ^{\prime} } = \begin{bmatrix}[r]
      -7 &-15   \\
       6&12   \\
    \end{bmatrix} \begin{bmatrix}[r]
       3 \\
       -2 \\
    \end{bmatrix}=\begin{bmatrix}[r]
       9 \\
       -6 \\
    \end{bmatrix} = [\mathbf{v}]_{\mathcal{B} ^{\prime} }
  \]
\end{solution}
\begin{exercise}
  Prove that if A and B are similar matrices, then $|A| =|B|$.
\end{exercise}
\begin{proof}
  Let the relation \(\sim \) denote similar matrices and let \(A,B\in M_{n\times n}(\mathbb{F}) \) have \(A\sim B\). Thus,
  \[
    \exists P\in \operatorname{GL}_n(\mathbb{F}) : A = P ^{-1} BP
  \]
  by definition of similar matrices. Define \(P\) to be such a matrix for which the above condition holds, and observe that
  \[
    A = P ^{-1} BP \Longrightarrow \det (A)=\det \left( P ^{-1} BP \right)
  \]
  From repeated applications of one of the more important theorems in Chapter 3,
  \[
    \det \left( P ^{-1} BP \right) = \det \left( P ^{-1}  \right) \det (B)\det (P)
  \]
  From another one of the theorems, we know that for any \(P\),
  \[
    \det \left( P ^{-1}  \right) = \frac{1}{\det (P)}
  \]
  Therefore, we have
  \begin{align*}
     \det (A)&=\det \left( P ^{-1} BP \right)\\
    &=\det \left( P ^{-1}  \right) \det (B)\det (P)\\
    &=\frac{\det (P)}{\det (P)}\det (B)\\
    &=\det (B)
  \end{align*}
  Since we define \(P\) such that \(P ^{-1} \) exists, \(\det (P)\neq 0\), and the above statements hold. Thus, \(\det (A)=\det (B)\), and therefore
  \[
    A \sim B \Longrightarrow \det (A) = \det (B)
  \]
\end{proof}
\begin{exercise}
  Prove that if matrix $A$ is diagonalizable, then $A^T$ is diagonalizable.
\end{exercise}
\begin{proof}
  For the sake of simplicity, let
  \[
    \mathfrak{D}_n (\mathbb{F})\coloneqq \left\{ A\in M_{n\times n}(\mathbb{F}) \mid A\text{ is diagonal} \right\}
  \]
  Let \(A\in M_{n\times n}(\mathbb{C})\). We say that \(A\) is diagonalizable if there exists some matrix \(P\in \operatorname{GL}_n(\mathbb{C}) \) such that
  \[
   P ^{-1}  AP \in \mathfrak{D}_n (\mathbb{C}) 
  \]
  where \(D\) is some diagonal matrix. In other words, \(A\) is similar to some diagonal matrix \(D\). Let \(P\) be a matrix where this is true, and define \(D\in \mathfrak{D}_n(\mathbb{C}) \) as
  \[
    D\coloneqq  P ^{-1}  AP
  \]
  It follows directly that
  \begin{align*}
    &\qquad A = PDP ^{-1} \\
    &\Longrightarrow A^{T} = \left( PDP ^{-1}  \right)^T
  \end{align*}
  Recall that for matrices \(X,Y\), \((XY)^T = Y^T X^T\). Thus,
  \begin{align*}
    A^{T} &= \left( PDP ^{-1}  \right)^T\\
    &= \left( P \left( DP ^{-1}  \right)  \right)^T\\
    &= \left( DP ^{-1}  \right)^T P^T\\
    &=\left( P ^{-1}  \right)^T D^T P^T
  \end{align*}
  Recall that for some \(X\in \operatorname{GL}_n(\mathbb{F}) \), \(\left( X ^{-1}  \right)^T =\left( X^T \right)^{-1}   \). Thus,
  \[
    A^T = \left( P ^{-1}  \right)^T D^T P^T = \left( P^T \right) ^{-1} D^T P^T
  \]
  Trivially, the transpose of a diagonal matrix is diagonal, so \(D^T \in \mathfrak{D}_n(\mathbb{C}) \). Furthermore, since \(P^T\) and \(\left( P^T \right)^{-1}  \) are inverses, they must be invertible matrices. Define \(C \coloneqq  P^T\), and as seen previously, \(C\in \operatorname{GL}_n(\mathbb{C}) \). We have 
  \[
    A^T = C^{-1} D^TC
  \]
  By definition of similar matrices, \(A^T\) is similar to \(D^T\), and therefore \(A^T\) is diagonalizable.
\end{proof}
\begin{exercise}
  Let $A=\begin{bmatrix*}[r] 2&-2&3\\0&3&-2\\0&-1&2 \end{bmatrix*}$.
\begin{enumerate}
\item Find the characteristic polynomial of $A$.
\item Find the eigenvalues of $A$. 
\item Find the corresponding eigenvectors for the eigenvalues found in part (b). 
\item Is the matrix A diagonalizable?  Why or why not? \\
\end{enumerate}
\end{exercise}
\begin{solution}
  Let \(\lambda \) be an eigenvalue for \(A\) with eigenvector \(\mathbf{x}\).
  \[
    \det \left( \lambda I - A \right) = \begin{vmatrix}[r]
      \lambda -2 &2  &-3   \\
       0&\lambda -3  &2   \\
       0&1  &\lambda -2   \\
    \end{vmatrix} = (\lambda -2)((\lambda -3)(\lambda -2)-2)=0
  \]
  \[
    \Longrightarrow (\lambda -2)\left( \lambda ^2 - 5\lambda +4\right) = \lambda ^3-7\lambda ^2+14\lambda+8=0
  \]
  is the characteristic polynomial of \(A\), which has three roots in \(\mathbb{C}\). Starting from the form \(0=(\lambda -2)\left( \lambda ^2 - 5\lambda +4\right)\) and factoring fully, we obtain
  \[
    0=(\lambda -2)(\lambda -4)(\lambda -1)
  \]
  Since the characteristic polynomial factors into three distinct linear factors over \(\mathbb{Q}\), we have three real-valued eigenvalues \(\lambda \in \{ 1,2,4 \} \). Take \(\mathbf{x}=(x,y,z)\) and let
  \[
    (\lambda I-A)\mathbf{x}=\mathbf{0}
  \]
  For \(\lambda =1\), it follows that
  \[
    \begin{bmatrix}[r]
      -1 &2  &-3   \\
       0&-2  &2   \\
       0&1  &-1   \\
    \end{bmatrix} \begin{bmatrix}[r]
       x \\
       y \\
       z \\
    \end{bmatrix} = \begin{bmatrix}[r]
       0 \\
       0 \\
       0 \\
    \end{bmatrix}\\
  \]
  We give the representation
  \[
    \begin{bmatrix}[r]
      -1 &2  &-3  &0   \\
       0&-2  &2  &0   \\
       0&1  &-1  &0   \\
    \end{bmatrix}\xlongrightarrow{rref} \begin{bmatrix}[r]
      1 &0  &1  &0   \\
       0&1  &-1  &0   \\
       0&0  &0  &0   \\
    \end{bmatrix}
  \]
  Thus, a basis for the eigenspace of \(\lambda =1\) is \(E_1 = \{ (-1,1,1) \} \). Similarly, for \(\lambda =2\),
  \[
    \begin{bmatrix}[r]
      0 &2  &-3  &0   \\
       0&-1  &2  &0   \\
       0&1  &0  &0   \\
    \end{bmatrix}\xlongrightarrow{rref} \begin{bmatrix}[r]
      0 &1  &0  & 0  \\
       0&0  &1  &0   \\
       0&0  &0  &0   \\
    \end{bmatrix}
  \]
  implies that \(E_2 = \{ (1,0,0) \} \). Finally, for \(\lambda =4\),
  \[
    \begin{bmatrix}[r]
      1 &2  &-3  &0   \\
       0&0  &2  &0   \\
       0&1  &1  &0   \\
    \end{bmatrix}\xlongrightarrow{rref} \begin{bmatrix}[r]
       1&0  &0  &0   \\
       0&1  &0  &0   \\
       0&0  &1  &0   \\
    \end{bmatrix}
  \]
  only the trivial solution exists. Thus, \(E_3 = \varnothing \) since the eigenspace is the trivial vector space. Notice that every eigenvector of \(A\) is linearly independent of the others. Therefore, \(A\) is diagonalizable.
\end{solution}
\begin{exercise}
  Let $T: P_1\rightarrow  P_1$ be a linear transformation such that
  \[T(a+ bx)=a+(a+2b)x\]
	If possible, find a basis $B$ for $P_1$ such that  the matrix for $T$ relative for $B$ is diagonal.
\end{exercise}
\begin{solution}
  For convenience, we represent \(P_1\) with \(\mathbb{R}^2\), since they are isomorphic. We give the matrix \(A\) for which \(T(\mathbf{x})=A \mathbf{x}\):
  \[
    A= \begin{bmatrix}[r]
       1&0   \\
       1& 2  \\
    \end{bmatrix}
  \]
  Trivially, this matrix has eigenvalues \(1\) and \(2\). This matrix will be diagonalizable. Now consider the reduced row-echelong form of the matrices \(\lambda I - A\).
  \[
    \begin{bmatrix}[r]
      0 &0  &0   \\
       -1&-1  &0   \\
    \end{bmatrix}\xlongrightarrow{rref} \begin{bmatrix}[r]
      1 &1&0   \\
       0& 0&0  \\
    \end{bmatrix}
  \]
  Thus, the eigenvector \((1,-1)\) forms a basis for the eigenspace given by \(\lambda =1\).
  \[
    \begin{bmatrix}[r]
      1 &0  &0   \\
       -1&0  &0   \\
    \end{bmatrix} \xlongrightarrow{rref} \begin{bmatrix}[r]
      1 &0&0   \\
       0&0&0   \\
    \end{bmatrix}
  \]
  Thus, the eigenvector \((0,1)\) forms a basis for the eigenspace given by \(\lambda =2\). Taking these eigenvectors, define
  \[
    P\coloneqq \begin{bmatrix}[r]
      1 & 0  \\
      -1 & 1  \\
    \end{bmatrix}
  \]
  which has inverse
  \[
    P ^{-1} = \begin{bmatrix}[r]
      1 &0   \\
       1&1   \\
    \end{bmatrix}
  \]
  Consider the matrix
  \[
    D\coloneqq P ^{-1} AP = \begin{bmatrix}[r]
      1 &0   \\
       1&1   \\
    \end{bmatrix}\begin{bmatrix}[r]
      1&0   \\
      1& 2  \\
   \end{bmatrix}\begin{bmatrix}[r]
      1 & 0  \\
      -1 & 1  \\
    \end{bmatrix}= \begin{bmatrix}[r]
      1 &0   \\
       0&  2 \\
    \end{bmatrix}
  \]
  which is diagonal. Thus, \(D \) is a diagonal matrix giving the transformation \(T\) on the basis induced by \(P ^{-1} \), which can be realized via the equivalent statement
  \[
    A = PDP ^{-1} 
  \]
  This statement makes it obvious that \(D\) represents the map \(T\) on the basis induced by \(P ^{-1} \). To solve the problem, we must identify this basis and give its representation in \(P_1\). First, let \(\mathbf{v}_1,\mathbf{v}_2\in \mathcal{B}\), and take
  \[
    \mathbf{v}_1 =\begin{bmatrix}[r]
      1 &0   \\
       1&1   \\
    \end{bmatrix} \begin{bmatrix}[r]
       1 \\
       0 \\
    \end{bmatrix} = \begin{bmatrix}[r]
       1 \\
       1 \\
    \end{bmatrix}
  \]
  and
  \[
    \mathbf{v}_2 =\begin{bmatrix}[r]
      1 & 0  \\
       1&1   \\
    \end{bmatrix} \begin{bmatrix}[r]
       0 \\
       1 \\
    \end{bmatrix} = \begin{bmatrix}[r]
       0 \\
       1 \\
    \end{bmatrix}
  \]
  Converting back to \(P_1\) we say that the basis set
  \[
    \mathcal{B}=\left\{ \begin{bmatrix}[r]
       1 \\
       x \\
    \end{bmatrix}, \begin{bmatrix}[r]
       0 \\
       x \\
    \end{bmatrix} \right\} 
  \]
  gives the transformation \(T\) as a diagonal matrix.
\end{solution}
\end{document}